\documentclass[11pt,a4paper]{article}
\usepackage[T1]{fontenc}
\usepackage[utf8]{inputenc}
\usepackage{lmodern}
\usepackage[margin=25mm]{geometry}
\usepackage{amsmath,amssymb,amsthm,mathtools}
\usepackage{booktabs,array}
\usepackage[hidelinks]{hyperref}
\usepackage{microtype}
\setlength{\parskip}{0.35em}
\setlength{\parindent}{1.2em}
\newtheorem{proposition}{Proposition}
\newtheorem{lemma}{Lemma}
\newcommand{\R}{\mathbb R}
\newcommand{\diag}{\operatorname{diag}}
\newcommand{\norm}[1]{\left\lVert#1\right\rVert}
\title{Adam in a Fixed Hessian Eigenbasis\\[0.35em]
\large AISC5020: Module 5 Method and Verification Note}
\author{HOU Kaiwen}
\date{15 October 2026}
\begin{document}
\maketitle
\begin{abstract}
This module studies a controlled intervention for coordinate dependence: apply
Adam in a fixed orthonormal eigenbasis of a positive-definite quadratic, then
map each update back to the original coordinates. We prove exact trajectory
correspondence with Adam on the diagonal problem, including positive numerical
stabilization and bias correction. We distinguish the common-basis statement
from the stronger, generally invalid claim that arbitrary eigenbases give the
same trajectory when eigenvalues are repeated. The implementation separates
basis preparation, residual diagnostics, optimization, and evaluation costs.
This note documents Module 5 only; it does not report a tuned comparison of all
team methods or claim completion of the group's main experiment.
\end{abstract}

\section{Question and controlled intervention}
Consider
\begin{equation}
 f(w)=\tfrac12 w^\top H w,\qquad
 H=VDV^\top,\qquad V^\top V=I,\qquad
 D=\diag(\lambda_1,\ldots,\lambda_p),\quad \lambda_i>0.
 \label{eq:problem}
\end{equation}
The eigenvalues determine the quadratic's spectral conditioning. They do not
specify the axes along which Adam accumulates squared gradient components.
Ordinary Adam is therefore an appropriate object for a coordinate experiment,
but changing coordinates and observing different convergence curves alone does
not identify which operation causes the difference.

Our intervention fixes the coordinates used for moment accumulation. With
$z=V^\top w$, the same objective becomes $\phi(z)=\tfrac12 z^\top Dz$.
At update $t$, let
\begin{equation}
 g_t=Hw_{t-1},\qquad h_t=V^\top g_t.
 \label{eq:gradient}
\end{equation}
We apply the bias-corrected Adam recurrence~\cite{adam} to $h_t$:
\begin{align}
 m_t&=\beta_1m_{t-1}+(1-\beta_1)h_t,
 &v_t&=\beta_2v_{t-1}+(1-\beta_2)(h_t\odot h_t),\label{eq:moments}\\
 \widehat m_t&=\frac{m_t}{1-\beta_1^t},
 &\widehat v_t&=\frac{v_t}{1-\beta_2^t},\label{eq:correction}\\
 d_t&=\frac{\widehat m_t}{\sqrt{\widehat v_t}+\epsilon\mathbf1},
 &w_t&=w_{t-1}-\eta_tVd_t.\label{eq:update}
\end{align}
The square, square root, and quotient of equal-length vectors are
componentwise; $Vd_t$ is a matrix--vector product. We use $m_0=v_0=0$,
$\beta_1,\beta_2\in[0,1)$, a common scalar
$\epsilon>0$ outside the square root, and shared step sizes $\eta_t>0$.
The moments live in eigen-coordinates even though $w_t$ is stored in ambient
coordinates. Transforming only the final direction while retaining ambient
second moments would define a different method.

\section{Exact correspondence and its scope}
\begin{proposition}[P4: correspondence in a specified common basis]
\label{prop:p4}
Run \eqref{eq:gradient}--\eqref{eq:update} on $H=VDV^\top$ from $w_0=Vz_0$.
Run ordinary Adam on $\phi(z)=\tfrac12 z^\top Dz$ from $z_0$, with the same
parameters, step sizes, zero initial moments, and update indexing. In exact
arithmetic, for every update $t$,
\begin{equation}
 w_t=Vz_t,\qquad f(w_t)=\phi(z_t).
 \label{eq:correspondence}
\end{equation}
The result allows repeated eigenvalues and positive $\epsilon$.
\end{proposition}
\begin{proof}
The initial points correspond by construction. Suppose $w_{t-1}=Vz_{t-1}$ and
the two pairs of moment vectors agree. Orthogonality and
\eqref{eq:problem} give
\[
 h_t=V^\top HVz_{t-1}=Dz_{t-1}=\nabla\phi(z_{t-1}).
\]
Thus \eqref{eq:moments} produces identical first and second moments in the two
runs. Their scalar bias corrections and positive scalar stabilizers also agree,
so their directions $d_t$ agree, including when a gradient component is zero.
Consequently,
\[
 w_t=Vz_{t-1}-\eta_tVd_t
     =V(z_{t-1}-\eta_td_t)=Vz_t.
\]
Finally, $f(Vz_t)=\tfrac12z_t^\top V^\top HVz_t=\phi(z_t)$.
Induction proves both claims.
\end{proof}

This is a trajectory identity, not a convergence-rate theorem. It gives no
guarantee that the chosen step size converges, that Adam outperforms GD, or that
the additional basis information is inexpensive. It also requires a fixed
basis: changing the basis during a run without consistently reinterpreting the
moment history invalidates the argument.

\begin{lemma}[Signed-permutation ambiguity]
\label{lem:sign}
Let $R$ be a signed permutation matrix and replace $V$ by $V'=VR$.
With corresponding coordinate initialization, the two eigenbasis-Adam runs
have the same ambient iterates in exact arithmetic.
\end{lemma}
\begin{proof}
Write $P=|R|$, an unsigned permutation matrix. Coordinate gradients and first
moments transform as $h'_t=R^\top h_t$ and $m'_t=R^\top m_t$, whereas second
moments transform as $v'_t=P^\top v_t$. This follows inductively because
$(R^\top h)\odot(R^\top h)=P^\top(h\odot h)$.
Bias correction commutes with both transformations. Since $P^\top\mathbf1=
\mathbf1$ and the denominator is permuted without a sign, the normalized
direction obeys $d'_t=R^\top d_t$. Therefore
$V'd'_t=VRR^\top d_t=Vd_t$, so the physical updates coincide.
\end{proof}

For a simple spectrum, every exact orthonormal eigenbasis of the same $H$
differs by such a signed permutation. A basis recovered from $H$ alone therefore
has the same idealized physical trajectory; aligning its signs to the generating
basis is unnecessary. Numerically close eigenvalues require more care because
their eigenvectors can be sensitive to small perturbations. In particular,
mathematically distinct eigenvalues do not by themselves guarantee numerically
resolvable eigenvectors or agreement of long floating-point trajectories.

For repeated eigenvalues, an eigensolver may choose an arbitrary orthogonal
rotation within a repeated eigenspace. Adam's componentwise square does not
commute with a general rotation, so Lemma~\ref{lem:sign} does not apply.
For example, let $H=I_2$, $w_0=(1,0)^\top$, and let $R$ be a rotation through
$\pi/4$. Both $I_2$ and $R$ are valid eigenbases. At the first update their
ambient Adam directions are, respectively,
\[
 \frac{(1,0)^\top}{1+\epsilon}
 \quad\text{and}\quad
 \frac{(1,0)^\top}{1/\sqrt2+\epsilon},
\]
which differ for every $\epsilon>0$. Proposition~\ref{prop:p4} remains valid when both runs use one
specified common basis; arbitrary independent eigenbasis recovery is a
different claim.

\section{Implementation and verification protocol}
Basis preparation is explicit and occurs before optimization. The implementation
supports three sources. An \emph{oracle} basis is the generating $V$ and isolates
the algebraic mechanism. An \emph{eigh} basis is obtained from the Hessian alone
using a symmetric eigensolver; the generating rotation is not used to align its
columns. A \emph{supplied} basis permits controlled sign, permutation, and
degenerate-subspace tests. The selected basis is copied, validated, and held
fixed, and a prepared object may be reused only for a matching Hessian.

We distinguish basis quality from trajectory agreement. For a candidate basis
$B=[b_1,\ldots,b_p]$ and
$\widetilde D=\diag(b_1^\top Hb_1,\ldots,b_p^\top Hb_p)$, report
\begin{equation}
 r_{\rm orth}=\frac{\norm{B^\top B-I}_F}{\sqrt p},\qquad
 r_{\rm diag}=\frac{\norm{HB-B\widetilde D}_F}{\norm H_F}.
 \label{eq:residuals}
\end{equation}
Positive eigenvalues, finite arrays, dimensions, symmetry, and residual
tolerances are checked separately. Numerical eigenvalue clusters are recorded
as a warning about basis ambiguity; a tolerance-based cluster is not proof of
an exactly repeated eigenvalue. The diagnostics also compare the minimum
eigengap with a conservative Hessian-scale allowance based on residuals,
roundoff, and the accepted validation tolerance. Stored eigenvalues may be
distinct while the eigenvector-uncertainty flag is true. This indicator is not
a certified forward-error bound for a particular eigensolver.

For paired trajectories, retain both the ambient discrepancy
\begin{equation}
 e_t=\frac{\norm{w_t-Vz_t}_2}{\max(1,\norm{w_0}_2)}
 \label{eq:trajectoryerror}
\end{equation}
and the absolute loss difference. A fixed initial scale avoids a misleading
relative error whose denominator vanishes near the optimizer. Residuals in
\eqref{eq:residuals} answer whether $B$ approximately diagonalizes $H$;
\eqref{eq:trajectoryerror} answers a different, dynamical question. Neither
should be reported as the other.

The verification suite covers the following distinct claims:
\begin{enumerate}
 \item Compare every recorded iterate for a shared generating basis, with both
 distinct and repeated spectra.
 \item Recover a basis from $H$ alone for a distinct spectrum and retain the
 eigensolver's signs and ordering.
 \item Supply a signed and permuted basis to check
 Lemma~\ref{lem:sign}; separately supply an admissible rotation within a repeated
 eigenspace to expose the limitation.
 \item Reject invalid bases and basis reuse with a different Hessian, and
 count preparation and transformation work explicitly.
\end{enumerate}
Each numerical record should preserve the spectrum, initial point, seed,
parameters, horizon, code version, and error definition. Module verification
uses fixed settings to test the mechanism. It does not replace the team's
independent tuning and statistical protocol.

\subsection{Observed structural verification}
The saved run of 15 October 2026 (08:36:01 UTC) uses seed 502015, 256 updates,
$\eta=0.01$, $(\beta_1,\beta_2)=(0.9,0.999)$, and $\epsilon=10^{-8}$.
There are 84 two-dimensional cases, 24 cases with $p=20$, 16 with $p=100$,
and two declared scope controls, for 126 cases in total. The main spectra use
$\kappa\in\{10,10^2,10^3,10^4\}$; the two-dimensional angles range from
$0^\circ$ to $90^\circ$ in $15^\circ$ increments. Each case has one diagonal
reference and three runs on the rotated problem: ordinary Adam, oracle-basis
Adam, and recovered-basis Adam. Thus 504 trajectories are saved, including
378 rotated-problem trajectories.

All 250 required correspondence checks passed the preset normalized position
tolerance $5\times10^{-8}$; Table~\ref{tab:errors} reports the observed maxima.
The recovered-basis requirement applies to the 124 separated-spectrum cases,
whereas the oracle common-basis requirement includes both controls. These
settings test an identity rather than tune the methods for minimum loss.

\begin{table}[htbp]
 \centering
 \caption{Maximum normalized path discrepancy over required checks.}
 \label{tab:errors}
 \begin{tabular}{lrr}
 \toprule
 Method and scope & Cases & Maximum $e_t$ \\
 \midrule
 Oracle basis, all cases & 126 & $1.535\times10^{-12}$ \\
 Recovered basis, separated spectra & 124 & $6.486\times10^{-13}$ \\
 \bottomrule
 \end{tabular}
\end{table}

The controls make the scope visible rather than silently omitting difficult
cases. For the repeated identity problem, the recovered-basis discrepancy is
$0.1819$, while the oracle common-basis discrepancy is $2.483\times10^{-16}$.
For the near-repeated spectrum
$(10^{-4},10^{-4}+10^{-11},1)$, the recovered-basis discrepancy is
$5.637\times10^{-7}$ despite a diagonalization residual
$4.885\times10^{-16}$. Its stored eigenvalues are distinct, but the numerical
uncertainty flag is true. The oracle discrepancy remains
$1.535\times10^{-12}$. This illustrates why a small backward residual and
distinct stored eigenvalues cannot alone establish forward trajectory agreement.

A separate checker recomputed all 504 saved trajectories from the saved
problems using explicit Adam recurrences, without the production Adam update
helper. It also independently regenerated the seeded case schedule and initial
points, recomputed losses, summary metrics, basis residuals and timing medians,
and checked operation counts, all 45 timing rows, and 19
source/configuration/script hashes. Its saved verification status
is passed. The environment was Python 3.12.14, NumPy 2.5.4, and Windows 11,
with BLAS thread environment variables set to one.

\section{Cost accounting and numerical interpretation}
For dense $p\times p$ matrices, recovering all eigenvectors costs
$O(p^3)$ arithmetic and $O(p^2)$ storage. Validation also contains dense
matrix--matrix products and belongs to setup. At each optimization step,
$V^\top g_t$ and $Vd_t$ add two dense matrix--vector products, each costing
$O(p^2)$, as well as $O(p)$ moment operations. The common quadratic gradient
$Hw$ already costs $O(p^2)$. Computing a loss for monitoring is additional
evaluation work and must not be silently counted as a free gradient call.

Our accounting separates basis setup, the full evaluation loop, and total
elapsed time. The evaluation loop includes gradients, coordinate transforms,
losses, checks, and trace bookkeeping; operation counts distinguish these
activities, but the implementation does not time each one separately.
Counts are more portable than timings on small matrices.
Reusing a prepared basis avoids a second eigendecomposition but
does not erase the first preparation cost; original setup and reuse metadata
must remain visible. Likewise, an oracle basis removes recovery from the timed
optimization loop by assumption, not by providing a generally free algorithm.

Table~\ref{tab:timings} gives local medians from five repetitions per dimension
and method, each with 256 updates. One untimed warm-up run per dimension and
method is excluded from the five measured repetitions. These are descriptive
timings, not evidence of a practical speedup. Setup and total medians are calculated
separately, so adding component medians need not reproduce a total median.

\begin{table}[htbp]
 \centering
 \caption{Local median time in milliseconds. Basis-method cells show
 setup / total; ordinary Adam has no basis setup.}
 \label{tab:timings}
 \begin{tabular}{rrrr}
 \toprule
 $p$ & Ordinary total & Oracle setup / total & Recovered setup / total \\
 \midrule
 2 & 12.508 & 0.304 / 14.995 & 0.248 / 11.474 \\
 20 & 11.635 & 0.435 / 13.160 & 0.468 / 12.723 \\
 100 & 11.281 & 1.852 / 18.791 & 2.341 / 15.621 \\
 \bottomrule
 \end{tabular}
\end{table}

The proof uses exact arithmetic, while implementations repeatedly multiply
floating-point matrices and accumulate moment histories. Small perturbations
can be amplified by normalization, especially when gradients and loss are
tiny. A long-horizon discrepancy after near-zero error is not a logical
counterexample to Proposition~\ref{prop:p4}. It is still relevant empirical
information: report its onset, magnitude, objective scale, and settings rather
than hiding it or treating every late difference as a failed eigensolver.

Conversely, a small decomposition residual alone does not establish a uniformly
small trajectory discrepancy for an arbitrary horizon. This report makes no
forward-error bound or stability theorem of that form. Interpretation should
combine algebra, independent implementation checks, basis residuals, and the
observed trajectories.

The reproducible artifacts are \path{outputs/part5/summary.json},
\path{correspondence.csv}, \path{basis_diagnostics.json}, \path{timing.csv},
the saved raw trajectories, and \path{artifact_verification.json}.
The same output directory contains \path{correspondence.pdf} and
\path{cost.pdf}; this standalone source has no external figure dependency.

\section{Integration and conclusion}
The intervention gives the group a precise mechanism test. If ordinary Adam
changes under a paired rotation while fixed-eigenbasis Adam tracks the diagonal
reference over the same regime, the result supports the role of coordinatewise
moment accumulation. That conditional interpretation must be based on actual
paired outputs; it does not follow from a planned figure or from the theorem
alone. Learning-rate selection, rotation sampling, uncertainty estimates, and
comparison with preconditioning remain the responsibility of the corresponding
team modules.

For integration, Module 5 contributes a fixed-basis implementation, a proof
with explicit assumptions, diagnostics that separate recovery from dynamics,
and cost metadata. The exact claim is deliberately limited: a specified common
eigenbasis restores correspondence with diagonal-coordinate Adam; a simple
spectrum makes signed-permutation ambiguity harmless; repeated eigenspaces
retain a genuine coordinate choice. These distinctions allow a stronger
mechanistic explanation without claiming universal rotation invariance or an
unmeasured practical speedup.

\begin{thebibliography}{1}
\bibitem{adam}
D. P. Kingma and J. Ba.
\newblock Adam: A Method for Stochastic Optimization.
\newblock \emph{International Conference on Learning Representations}, 2015.
\newblock \url{https://arxiv.org/abs/1412.6980}.
\end{thebibliography}
\end{document}
